% Auto-generated by Code2Latex
% Requires: longtable, booktabs packages

\begin{longtable}{p{\textwidth}}
\caption{Tools for single (Level 1)} \\
\endfirsthead

\multicolumn{1}{c}{\tablename\ \thetable{} -- continued from previous page} \\
\endhead

\multicolumn{1}{r}{Continued on next page} \\
\endfoot

\endlastfoot

\textbf{\texttt{add\_adsorbate\_to\_slab\_text}} \\
\begin{description}
    \item[\textbf{Description:}] Place an adsorbate molecule on a surface slab at a specified adsorption site.

This tool combines a surface slab with an adsorbate molecule by placing the adsorbate at a specific binding site on the surface.
    It handles the geometric placement of the molecule at the correct height above the surface and ensures proper structural integration.
    This is essential for creating realistic surface-adsorbate systems for computational studies of catalysis, adsorption energetics, and surface reactivity.
    \item[\textbf{Arguments:}]
    \begin{description}
        \item \texttt{slab\_cif} (str): CIF string of the surface slab structure.
A properly formatted CIF string containing the surface slab structure on which the adsorbate will be placed. This should be a two-dimensional periodic structure with a well-defined surface and vacuum region. The slab provides the substrate for molecular adsorption.
        \item \texttt{adsorbate\_cif} (str): CIF string of the adsorbate molecule structure.
A CIF or XYZ formatted string containing the molecular structure of the adsorbate to be placed on the surface. This can be a small molecule like CO2, H2O, or more complex organic molecules. The tool will attempt to parse both CIF and XYZ formats automatically.
        \item \texttt{height} (float): Height in Angstroms above the surface for adsorbate placement. Defaults to 2.0.
The vertical distance above the surface at which the adsorbate will be placed. This parameter controls the initial separation between the adsorbate and the surface atoms. Typical values range from 1.5 to 3.0 Å depending on the molecular size and expected binding interaction.
        \item \texttt{site} (list[float] | None): Optional fractional coordinates for adsorbate placement.
A list of three floating-point numbers representing the fractional coordinates [x, y, z] where the adsorbate should be placed If not provided, the tool will automatically select the first available top site. These coordinates should typically come from choose\_adsorption\_site\_text output.
    \end{description}
    \item[\textbf{Returns:}] CIF string of the combined surface-adsorbate structure.

A properly formatted CIF string containing the combined structure with the adsorbate placed on the surface at the specified position and height.
             This structure includes both the original slab atoms and the adsorbate atoms, properly integrated into a single periodic structure.
\end{description} \\
\midrule
\textbf{\texttt{cat\_files}} \\
\begin{description}
    \item[\textbf{Description:}] Concatenate and display contents of one or more files
    \item[\textbf{Arguments:}]
    \begin{description}
        \item \texttt{paths} (list[str]): List of file paths to concatenate
        \item \texttt{separator} (str): Separator between file contents
    \end{description}
\end{description} \\
\midrule
\textbf{\texttt{choose\_adsorption\_site\_text}} \\
\begin{description}
    \item[\textbf{Description:}] Select a specific adsorption site from classified sites by type and index.

This tool selects one specific adsorption site from a collection of classified sites based on the site type (ontop, bridge, hollow) and index within that type.
This selection is crucial for systematic studies of different binding environments and their effects on adsorption energetics.
    \item[\textbf{Arguments:}]
    \begin{description}
        \item \texttt{adsorption\_sites\_json} (str): JSON string mapping site types to lists of fractional coordinates.
A JSON-formatted string containing a dictionary where keys are site types (e.g., "ontop", "bridge", "hollow") and values are lists of fractional coordinates. This should be the output from get\_adsorption\_sites\_text. Each coordinate is a list of three numbers representing position within the unit cell.
        \item \texttt{site\_type} (str): Type of adsorption site to select.
The type of binding site to select from the available options. Common types include "ontop" (above surface atoms), "bridge" (between two atoms), and  "hollow" (in multi-atom depressions). The type must exist in the JSON dictionary and determines the coordination environment of the selected site.
        \item \texttt{index} (int): Index of the site within the specified type. Defaults to 0.
The numerical index of the site to select from the list of sites of the specified type. Index 0 selects the first site, index 1 the second, and so on. The index must be within the range of available sites for the specified type.
    \end{description}
    \item[\textbf{Returns:}] Fractional coordinates of the selected adsorption site.

A list of three floating-point numbers representing the fractional coordinates [x, y, z] of the selected adsorption site within the unit cell.
                These coordinates can be used directly for adsorbate placement and represent the optimal binding position for the specified site type.
\end{description} \\
\midrule
\textbf{\texttt{choose\_slab\_text}} \\
\begin{description}
    \item[\textbf{Description:}] Select a specific slab from a JSON dictionary of enumerated slabs by index.

This tool selects one surface slab from a collection of enumerated slabs based on its index number.
It's designed to work with the output from enumerate\_slabs\_text, allowing users to choose a specific surface termination for further analysis.
This selection process is crucial for focusing on the most relevant or interesting surface termination for catalysis or adsorption studies.
    \item[\textbf{Arguments:}]
    \begin{description}
        \item \texttt{slabs\_json} (str): JSON string mapping slab keys to CIF strings.
A JSON-formatted string containing a dictionary where keys are slab identifiers (e.g., "slab\_0", "slab\_1") and values are the corresponding CIF strings. This should be the output from enumerate\_slabs\_text tool. The JSON structure must be valid and contain at least one slab entry.
        \item \texttt{index} (int): Index of the slab to select. Defaults to 0.
The numerical index of the slab to select from the JSON dictionary. This corresponds to the enumeration order from enumerate\_slabs\_text, where index 0 is the first slab, index 1 is the second, and so on. The tool will look for a key named "slab\_\{index\}" in the JSON dictionary.
    \end{description}
    \item[\textbf{Returns:}] CIF string for the selected slab.

A properly formatted CIF string containing the structure data for the selected slab.
         This includes atomic positions, lattice parameters, and all necessary crystallographic information.
         The CIF can be used directly with other structure analysis tools.
\end{description} \\
\midrule
\textbf{\texttt{copy\_file}} \\
\begin{description}
    \item[\textbf{Description:}] Copy a file from a source to destination
    \item[\textbf{Arguments:}]
    \begin{description}
        \item \texttt{source} (str): Source file path
        \item \texttt{destination} (str): Destination file path
    \end{description}
\end{description} \\
\midrule
\textbf{\texttt{enumerate\_slabs\_text}} \\
\begin{description}
    \item[\textbf{Description:}] Enumerate all possible surface slab terminations from a bulk structure and return as JSON.

This tool generates all possible surface terminations for a given bulk crystal structure along specified Miller indices.
    Unlike creating a single slab, this tool explores different ways to terminate the surface, which is crucial for materials with complex structures or multiple chemically distinct layers.
    Each termination represents a different surface chemistry and reactivity, making this tool essential for comprehensive surface studies.
    \item[\textbf{Arguments:}]
    \begin{description}
        \item \texttt{bulk\_cif} (str): Bulk crystal structure in CIF string format.
A properly formatted CIF string containing the bulk crystal structure from which surface slabs will be generated. This should be a three-dimensional periodic structure with well-defined atomic positions and lattice parameters. The structure will be analyzed to determine all possible surface terminations.
        \item \texttt{miller\_index} (tuple): Miller indices for surface orientation. Defaults to (1,1,1).
A tuple of three integers specifying the crystallographic plane along which all surface terminations will be generated. This determines the surface orientation but allows for different terminations along the same plane. Different Miller indices will produce different surface structures and properties.
        \item \texttt{min\_slab\_size} (float): Minimum slab thickness in Angstroms. Defaults to 12.
The minimum thickness of each slab in the direction perpendicular to the surface plane. This ensures that all generated slabs have sufficient bulk-like character while exposing different surface terminations. Affects both the structural accuracy and computational requirements.
        \item \texttt{min\_vacuum\_size} (float): Minimum vacuum layer thickness in Angstroms. Defaults to 5.
The minimum vacuum space above each surface to prevent interactions between periodic images. This parameter is applied to all generated slabs and is crucial for accurate surface calculations. Larger values reduce spurious interactions but increase computational cost.
    \end{description}
    \item[\textbf{Returns:}] JSON string mapping slab indices to their CIF representations.

A JSON-formatted string containing a dictionary where keys are slab identifiers (e.g., "slab\_0", "slab\_1") and values are the corresponding CIF strings for each surface termination.
             This format allows easy selection and comparison of different surface terminations.
\end{description} \\
\midrule
\textbf{\texttt{file\_info}} \\
\begin{description}
    \item[\textbf{Description:}] Get information about a file or directory
    \item[\textbf{Arguments:}]
    \begin{description}
        \item \texttt{path} (str): Path to the file or directory
    \end{description}
\end{description} \\
\midrule
\textbf{\texttt{get\_adsorption\_sites\_text}} \\
\begin{description}
    \item[\textbf{Description:}] Identify and classify all possible adsorption sites on a surface slab.

This tool analyzes a surface slab structure to identify and classify potential adsorption sites where molecules can bind.
    It uses geometric and chemical analysis to determine different types of binding sites such as top sites (above surface atoms), bridge sites (between two atoms), and hollow sites (in multi-atom depressions).
    This analysis is fundamental for understanding surface reactivity and designing catalysts.
    \item[\textbf{Arguments:}]
    \begin{description}
        \item \texttt{slab\_cif} (str): CIF string of the surface slab structure.
A properly formatted CIF string containing the surface slab structure with atomic positions, lattice parameters, and surface geometry. This should be a two-dimensional periodic structure with a well-defined surface and vacuum region. The structure is analyzed to identify potential adsorption sites.
    \end{description}
    \item[\textbf{Returns:}] JSON string containing classified adsorption sites with fractional coordinates.

A JSON-formatted string containing a dictionary where keys are site types (e.g., "top", "bridge", "hollow") and values are lists of fractional coordinates for each site of that type.
             Each coordinate is a list of three numbers [x, y, z] representing the fractional position within the unit cell.
\end{description} \\
\midrule
\textbf{\texttt{get\_structure\_from\_mp\_text}} \\
\begin{description}
    \item[\textbf{Description:}] Retrieve a pymatgen structure from Materials Project using its API and return CIF content as text.

This tool connects to the Materials Project database to download crystal structure data for a given material ID.
    It retrieves the structure object and converts it to CIF (Crystallographic Information File) format, which is the standard format for storing crystal structure information.
    CIF is then returned as string
    \item[\textbf{Arguments:}]
    \begin{description}
        \item \texttt{mp\_id} (str): Materials Project identifier string.
The unique identifier used by Materials Project to catalog materials. Should be in the format "mp-XXXXX" where XXXXX is a numerical ID. This ID corresponds to a specific material entry in the Materials Project database.
    \end{description}
    \item[\textbf{Returns:}] CIF content string containing the crystal structure data.

A properly formatted CIF (Crystallographic Information File) string containing all necessary information about the crystal structure including lattice parameters, atomic positions, space group, and symmetry operations.
             This format is widely compatible with crystallographic software and other structure analysis tools.
\end{description} \\
\midrule
\textbf{\texttt{list\_files}} \\
\begin{description}
    \item[\textbf{Description:}] List files in a directory
    \item[\textbf{Arguments:}]
    \begin{description}
        \item \texttt{path} (str): Path to the directory
        \item \texttt{recursive} (bool): Whether to list files recursively
    \end{description}
\end{description} \\
\midrule
\textbf{\texttt{read\_file}} \\
\begin{description}
    \item[\textbf{Description:}] Read the contents of a file into a string
    \item[\textbf{Arguments:}]
    \begin{description}
        \item \texttt{path} (str): Path to the file to read
    \end{description}
\end{description} \\
\midrule
\textbf{\texttt{write\_file}} \\
\begin{description}
    \item[\textbf{Description:}] Write content to a file
    \item[\textbf{Arguments:}]
    \begin{description}
        \item \texttt{path} (str): Path to the file to write
        \item \texttt{content} (str): Content to write to the file
    \end{description}
\end{description} \\
\end{longtable}
